\documentclass[10pt,a4paper]{article}

\usepackage[margin=1.05in]{geometry}
\usepackage{setspace}
\usepackage{amsmath,amssymb,amsthm}
\usepackage{booktabs}
\usepackage{longtable}
\usepackage{array}
\usepackage{caption}
\usepackage{enumitem}
\usepackage{titlesec}
\usepackage{microtype}
\usepackage{hyperref}
\usepackage{fancyhdr}

\hypersetup{colorlinks=true,linkcolor=black,urlcolor=black,citecolor=black}
\captionsetup{font=small,labelfont=bf,skip=4pt,justification=raggedright}

\titleformat{\section}{\normalsize\bfseries}{}{0em}{}
\titleformat{\subsection}{\small\bfseries}{}{0em}{}
\titleformat{\subsubsection}{\small\itshape}{}{0em}{}
\titlespacing*{\section}{0pt}{1.35em}{0.5em}
\titlespacing*{\subsection}{0pt}{0.95em}{0.35em}

\theoremstyle{plain}
\newtheorem{proposition}{Proposition}
\newtheorem{lemma}{Lemma}
\newtheorem{theorem}{Theorem}
\newtheorem{corollary}{Corollary}
\theoremstyle{definition}
\newtheorem{definition}{Definition}
\theoremstyle{remark}
\newtheorem{observation}{Observation}
\newtheorem{remark}{Remark}

\setstretch{1.14}
\setlist[itemize]{leftmargin=1.15em,itemsep=0.1em,topsep=0.2em}
\setlist[enumerate]{leftmargin=1.15em,itemsep=0.1em,topsep=0.2em}

\pagestyle{fancy}
\fancyhf{}
\fancyhead[L]{\small\itshape Elements of Position}
\fancyhead[R]{\small\itshape \ The Geometric Kernel}
\fancyfoot[C]{\small\thepage}
\renewcommand{\headrulewidth}{0pt}

\begin{document}

\begin{center}
\vspace*{2.2em}
{\large Elements of Position}\\[1.6em]
{\LARGE\bfseries\ The Geometric Kernel}\\[2.2em]
{\normalsize Justin Erholtz}\\[0.35em]
{\small 18 September 2026}
\end{center}

\vspace{2.4em}

\begin{center}
{\large What the walk measures}
\end{center}

\vspace{1.8em}

\noindent
Paper~I named the closed unit and the pair.
This paper is the walk that produces those measurements from two operators and a count.
Nothing is loaded into the operators.
What is named after the fact --- a half-turn length, a leftover of two readings, a ratio of chord to path --- is read off the tables the walk writes.

The source is \texttt{Elements\_of\_Position\_Geometric\_Kernel.py}.
The tables are \texttt{eop\_crossings.csv}, \texttt{eop\_pairs.csv}, \texttt{eop\_summary.csv}, from a run with shear \(h=0.002\) and \(400\) horizon hits.
Those numbers are a sheet size and a depth of count.
The identities do not belong to that sheet.

\vspace{1.4em}

\noindent
\textit{A note.}
Claims that can be proved from the operators are proved as identities.
Claims that are counts are reported from the tables.
The two are kept apart.

\newpage
\tableofcontents
\newpage

% =====================================================================
\section{The two operators}
% =====================================================================

Start at \((1,0)\).
A heading is a pair of coordinates whose length after the bind is \(1\).

\begin{definition}[Shear]
For a step size \(h>0\),
\[
T_h(x,y)=(x-hy,\,y+hx).
\]
\end{definition}

\begin{definition}[Bind]
\[
N(x,y)=\frac{1}{\sqrt{x^2+y^2}}(x,y).
\]
\end{definition}

\begin{definition}[Step]
\[
\mathrm{step}(x,y)=N\bigl(T_h(x,y)\bigr).
\]
\end{definition}

If the last state already had length \(1\), the un-bound length after the shear is an identity of the sheet:
\[
\bigl|T_h(x,y)\bigr|=\sqrt{1+h^2}.
\]
The bind leftover of one kick is therefore
\[
\sqrt{1+h^2}-1,
\]
constant in the step index.
On the recorded run this is \(1.999998\times 10^{-6}\).
The Euclidean distance between two successive bound headings is likewise a function of \(h\) only.
On the recorded run it is \(0.001999997\).
Those two numbers describe the grain of the walk.
They are not a position on the pair.

% =====================================================================
\section{The horizon count}
% =====================================================================

A hit is the heading crossing the line \(y=0\).
Down: \(y>0\) then \(y\le 0\).
Up: \(y<0\) then \(y\ge 0\).
Each hit is one half-turn of the heading.

\begin{definition}[First gap]
\(H_0\) is the step index of the first hit.
\end{definition}

On the recorded run \(H_0=1571\), side down, heading
\[
(-0.99999992,\,-0.00040316).
\]
The heading is on the left horizon to one kick of vertical.

\begin{definition}[Snap scale]
The \(m\)-th hit has snap scale \(k=m\).
\end{definition}

The frozen scale \(k_{\mathrm{frozen}}=j/H_0\) uses the first gap as a fixed ruler.
The two scales agree at \(m=1\).
They part when later gaps are not exactly \(H_0\).

\subsection{Gaps}

After the first hit the tables show only two gap values, \(1570\) and \(1571\).
Histogram on \(400\) hits:
\[
1571\text{ occurs }320\text{ times},\qquad
1570\text{ occurs }80\text{ times}.
\]
Mean gap:
\[
1570.8.
\]
The length a continuous turn of this shear would take is \(\pi/\mathrm{atan}(h)=1570.798421\).
The discrete walk cannot land on that length, so it pays with the two neighboring integers.
The mean tracks the continuous half-turn to three decimal places.

\begin{observation}
\(H_0\cdot h=3.142\) on this sheet.
That product is the half-turn read as steps times step size.
It is not inserted into \(T_h\) or \(N\).
\end{observation}

\subsection{Heading at the hits}

At every recorded hit the heading sits on a horizon.
Odd \(m\) (down): \(h_x\approx -1\), \(h_y\) in \((-0.002,0]\).
Even \(m\) (up): \(h_x\approx +1\), \(h_y\) in \([0,0.002)\).
The vertical coordinate never grows with \(m\).
Sampling at the cross keeps the heading on the wall.
Sampling at \(m\cdot H_0\) does not: that ruler stacks the first-crossing leftover and the heading walks off the axis.

% =====================================================================
\section{The pair}
% =====================================================================

At snap scale \(k=m\),
\[
r_{\mathrm{out}}=\frac k2,\qquad
r_{\mathrm{in}}=\frac2k.
\]

\begin{proposition}
\label{prop:product}
For every \(k>0\),
\[
r_{\mathrm{out}}\,r_{\mathrm{in}}=1.
\]
\end{proposition}

\begin{proof}
Immediate.
\end{proof}

\begin{proposition}
\label{prop:gm}
\[
\mathrm{GM}(r_{\mathrm{out}},r_{\mathrm{in}})=1.
\]
\end{proposition}

\begin{proof}
\(\sqrt{r_{\mathrm{out}}r_{\mathrm{in}}}=\sqrt{1}=1\).
\end{proof}

The crossings table records \(\mathrm{gm}=1\) on every row, to machine precision.
One row at \(m=374\) prints \(0.9999999999999999\).
That is the same identity written in floating point.

\begin{definition}[Leftover of the pair]
\[
\tau(k)=\mathrm{AM}-\mathrm{GM}
=\frac{k/2+2/k}{2}-1
=\frac k4+\frac1k-1.
\]
\end{definition}

\begin{proposition}
\label{prop:kiss}
\(\tau(k)=0\) if and only if \(k=2\) (for \(k>0\)).
\end{proposition}

\begin{proof}
\(\frac k4+\frac1k=1\) multiplies to \(k^2-4k+4=0\), hence \((k-2)^2=0\).
\end{proof}

The table: \(m=2\), \(\tau=0\), \(r_{\mathrm{out}}=r_{\mathrm{in}}=1\), side up, \(j=3142=2H_0\).
That is the kiss.
At \(m=1\), \(\tau=1/4\).
At \(m=4\), \(\tau=1/4\) again.

\begin{proposition}
\label{prop:fold}
\(\tau(k)=\tau(4/k)\) for every \(k>0\).
\end{proposition}

\begin{proof}
Substitute.
\end{proof}

The walls swap:
\[
r_{\mathrm{out}}(k)=r_{\mathrm{in}}(4/k),\qquad
r_{\mathrm{in}}(k)=r_{\mathrm{out}}(4/k).
\]
The pairs table records this at \(m=1\) with partner \(m=4\): leftover \(0.25\) on both sides, walls \((0.5,2)\) against \((2,0.5)\).
At \(m=2\) the state is its own partner.

The pairs file as written only keeps partners that round to an integer snap inside the run.
For \(k>4\) the partner \(4/k\) lies below \(1\) and is discarded by that rounding.
The identity \(\tau(k)=\tau(4/k)\) does not need the partner to be an integer snap.
It is an identity of the pair formula.

% =====================================================================
\section{Chord and path of one half-turn}
% =====================================================================

Place the outward wall on the heading:
\[
P_m=r_{\mathrm{out}}(m)\,(h_x,h_y).
\]
The straight line from \(P_{m-1}\) to \(P_m\) is the leg chord.
The summed micro-path between those hits is the leg path.
The table column \(\mathrm{leg\_chord\_over\_path}\) is their ratio.

\begin{proposition}
If the heading at hit \(m-1\) is the opposite of the heading at hit \(m\), and the radii are \((m-1)/2\) and \(m/2\), the chord is
\[
\frac{m-1}{2}+\frac m2=m-\tfrac12.
\]
\end{proposition}

\begin{proof}
Opposite unit headings \(u\) and \(-u\):
\[
\bigl|(m/2)(-u)-((m-1)/2)u\bigr|=\bigl(m/2+(m-1)/2\bigr)|u|=m-\tfrac12.
\]
\end{proof}

The table: \(m=2\) chord \(1.500000\); \(m=3\) chord \(2.500000\); \(m=4\) chord \(3.500000\).
The pattern holds.

The ratio of that chord to the summed path climbs toward \(2/\pi\).
Selected rows:

\begin{center}
\begin{tabular}{@{}rrr@{}}
\toprule
\(m\) & chord/path & \(2/\pi\) \\
\midrule
2 & \(0.62243\) & \(0.63662\) \\
4 & \(0.63402\) & \\
10 & \(0.63665\) & \\
20 & \(0.63689\) & \\
100 & \(0.63695\) & \\
400 & \(0.63654\) & \\
\bottomrule
\end{tabular}
\end{center}

The half-turn of a circle has diameter over semicircle equal to \(2/\pi\).
That is what the long baseline is reading.
The small oscillation at large \(m\) is the leftover angle of a discrete hit, not a second constant.

% =====================================================================
\section{The leftover integrated}
% =====================================================================

From the kiss outward the leftover of the pair is a region between the two walls and the axle \(\mathrm{GM}=1\).
Its integral in the snap scale is
\[
\int_2^N\tau(k)\,dk
=\Bigl[\frac{k^2}{8}+\ln k-k\Bigr]_2^N
=\frac{N^2}{8}+\ln N-N+\frac32-\ln 2.
\]

On the recorded run \(N=400\):
\[
\frac{400^2}{8}+\ln 400-400+\frac32-\ln 2=19606.798317.
\]
The summary table:
\[
\texttt{wedge\_from\_crossings}=19606.798317366545,
\]
\[
\texttt{wedge\_closed}=19606.79831736655,
\]
difference \(-3.6\times 10^{-12}\).
The script sums the same antiderivative from hit to hit.
The match is the identity, not a numerical accident.

When the leftover is instead accumulated as a rectangle sum along a frozen-\(H_0\) walk, the difference from the closed form at large \(N\) is
\[
\frac{\text{discrete}-\text{closed}}{N\,h}\approx\frac{1}{8\pi}.
\]
Check at an earlier frozen run, \(N=3600\), \(h=0.002\):
\[
\frac{0.28612}{3600\times 0.002}=0.03974,\qquad
\frac{1}{8\pi}=0.03979.
\]
That ratio is the staircase error on the \(k^2/8\) term once \(H_0\sim \pi/h\).
It is a property of summing on the frozen ruler.
The snap integral does not carry it.

% =====================================================================
\section{What the three files hold}
% =====================================================================

\subsection{\texttt{eop\_summary.csv}}

\begin{center}
\begin{tabular}{@{}ll@{}}
\toprule
key & value \\
\midrule
\(h\) & \(0.002\) \\
\(H_0\) & \(1571\) \\
half-turns & \(400\) \\
sheet leftover & \(1.999998\times 10^{-6}\) \\
heading chord & \(0.001999997\) \\
\(\pi/\mathrm{atan}(h)\) & \(1570.798421\) \\
\(H_0\cdot h\) & \(3.142\) \\
wedge (snap) & \(19606.798317\) \\
wedge (closed) & \(19606.798317\) \\
mean gap & \(1570.8\) \\
gap min / max & \(1570\) / \(1571\) \\
histogram & \(\{1571:320,\,1570:80\}\) \\
\bottomrule
\end{tabular}
\end{center}

Sheet: \(h\), leftover of one kick, heading chord, \(H_0\).
Object: product of the walls, leftover formula, kiss at \(2\), mean gap tracking the continuous half-turn, wedge identity.

\subsection{\texttt{eop\_crossings.csv}}

Four hundred rows, one per hit.
Forced on every row: \(\mathrm{gm}=1\), \(r_{\mathrm{out}}r_{\mathrm{in}}=1\), \(\tau=k/4+1/k-1\), heading on a horizon.
Side alternates down, up, down, up.
Frozen scale at \(m=400\) is \(399.949\), snap scale is \(400\).
The ruler \(j/H_0\) has slipped by five hundredths of a half-turn.
The heading has not.

\subsection{\texttt{eop\_pairs.csv}}

Integer partners of \(k\leftrightarrow 4/k\) that land inside the run.
The only clean integer pair in \(1\le k\le 400\) is \(1\leftrightarrow 4\), and the self-pair \(2\leftrightarrow 2\).
Rows that list partner \(m=1\) for \(k>4\) are the rounding rule, not the involution.
Read the involution from the formula and from \(m=1,2,4\).

% =====================================================================
\section{What stands}
% =====================================================================

The operators produce a heading that counts its own half-turns.
The count is two neighboring integers whose mean is the continuous half-turn of that shear.
Scale is the hit index.
The pair at that scale has product \(1\) and leftover \(\frac k4+\frac1k-1\), zero only at \(2\).
Chord over a half-turn of the outward wall, against the path of that same half-turn, reads \(2/\pi\).
The leftover integrated from the kiss is the elementary antiderivative of \(\tau\).
A frozen first-crossing ruler rotates the heading off the axis; a snap to the axis does not.

Named after the measure, not before it: \(H_0\cdot h\) as a half-turn length; \(2/\pi\) as diameter over semicircle; \(1/(8\pi)\) as the frozen-sum error.
None of those names enters \(T_h\) or \(N\).

What is not granted: a last decimal as closure of the heading; a completed run of the count; a leftover of the pair treated as a capacity or a field.

The first sentence of Paper~I is still the sentence.
The kernel is how the walk answers it.

\end{document}
